\documentclass[11pt,a4paper]{article}

% ── Packages ──────────────────────────────────────────────────────────────────
\usepackage[margin=2.5cm]{geometry}
\usepackage{amsmath,amssymb,amsthm}
\usepackage{booktabs}
\usepackage{graphicx}
\usepackage{subcaption}
\usepackage{hyperref}
\usepackage{xcolor}
\usepackage{enumitem}
\usepackage{multirow}
\usepackage{array}
\usepackage{rotating}
\usepackage{float}
\usepackage{microtype}
\usepackage{parskip}
\usepackage[backend=biber,style=authoryear,maxnames=2]{biblatex}

\hypersetup{
  colorlinks=true, linkcolor=blue!60!black,
  citecolor=green!50!black, urlcolor=blue!60!black
}

\definecolor{mblue}{RGB}{33,118,174}
\definecolor{mred}{RGB}{232,72,85}
\definecolor{mgray}{RGB}{100,100,100}

% ── Title ─────────────────────────────────────────────────────────────────────
\title{%
  \textbf{Membership Inference Attacks on\\
  Masked Diffusion Language Models}\\[6pt]
  \large Run~3 Full Report: Six MIMIR Domains\\[4pt]
  \normalsize White-Box Signal Extraction, SAMA Attack, and Ablation Study\\[6pt]
  \texttt{Model:}~\texttt{dllm-hub/Qwen3-0.6B-diffusion-mdlm-v0.1}
}
\author{Shailesh Kumar\\
  \small DA5001 — Privacy in AI\\
  \small \texttt{shailesh.mkvr@gmail.com}
}
\date{April 2026}

\begin{document}
\maketitle

% ── Abstract ──────────────────────────────────────────────────────────────────
\begin{abstract}
We present a comprehensive empirical study of membership inference attacks (MIA)
on a masked diffusion language model (MDLM), specifically
\texttt{Qwen3-0.6B-diffusion-mdlm-v0.1}, evaluated across six domains of the
MIMIR benchmark: ArXiv, GitHub, HackerNews, PubMed Central, Pile-CC, and Wikipedia.
Our pipeline fine-tunes the model on 1{,}000 member documents per domain and
evaluates three classes of attacks: black-box baselines (Loss, Zlib, Ratio),
the SAMA grey-box attack, and a white-box XGBoost/MLP classifier trained on
46 hand-crafted diffusion-trajectory signals.
The white-box classifier achieves AUC-ROC of $0.930$ on Pile-CC and $0.928$
on HackerNews, substantially outperforming SAMA ($0.885$ and $0.833$
respectively) and all black-box baselines.
A leave-one-group-out (LOSO) ablation over 13 signal groups shows that the
ELBO trajectory group is by far the most discriminative (mean LOSO AUC drop
$+0.135$), followed by predictive entropy ($+0.043$), with attention-based
features providing little additional signal.
All code, extracted features, and pre-computed ablation tensors are released.
\end{abstract}

\tableofcontents
\vspace{1em}
\hrule
\vspace{1em}

% ══════════════════════════════════════════════════════════════════════════════
\section{Introduction}
% ══════════════════════════════════════════════════════════════════════════════

Membership inference attacks (MIAs) ask whether a given data point was part of a
model's training set. While extensively studied for discriminative models and,
more recently, autoregressive language models, their applicability to
\emph{masked diffusion language models} (MDLMs) remains largely unexplored.

MDLMs define a forward process that gradually masks tokens at different rates
and train a transformer to predict the original tokens from their noised
counterparts. The training objective is a variational lower bound (ELBO) on
the data log-likelihood. Unlike autoregressive models, MDLMs do not produce a
natural per-token log-probability, making standard loss-based attacks less
direct. However, the diffusion trajectory itself --- how quickly the model's
loss decreases as the masking ratio is varied --- provides a rich
\emph{white-box} fingerprint of whether a document was memorised.

\textbf{This report} documents Run~3 of our MIA pipeline: the first large-scale
evaluation across all six MIMIR domains using a unified set of 46 diffusion-
trajectory features, the official SAMA attack, and a full ablation study that
decomposes which feature groups drive membership detection.

\subsection{Model}
\label{sec:model}

The target model is \texttt{dllm-hub/Qwen3-0.6B-diffusion-mdlm-v0.1}, a
Qwen3-0.6B transformer backbone adapted for masked diffusion pre-training.
Key properties:
\begin{itemize}[nosep]
  \item \textbf{Architecture}: Qwen3 transformer, 0.6B parameters.
  \item \textbf{Mask token ID}: 151{,}669 (Qwen3-specific).
  \item \textbf{Training objective}: MDLM-style ELBO; the forward process is
        linear $q(x_t | x_0) = \text{Mask}(x_0,\, t)$, $t \sim \mathcal{U}(\varepsilon, 1)$.
  \item \textbf{Inference}: Requires \texttt{trust\_remote\_code=True} and
        \texttt{Qwen2Tokenizer}; loading via \texttt{AutoModelForMaskedLM} with
        \texttt{attn\_implementation=``eager''} for attention-capture.
\end{itemize}

% ══════════════════════════════════════════════════════════════════════════════
\section{Dataset}
% ══════════════════════════════════════════════════════════════════════════════

\subsection{MIMIR Benchmark}
\label{sec:mimir}

We use the MIMIR benchmark \citep{duan2024mimir}, specifically the
\texttt{iamgroot42/mimir} Hugging Face dataset with split
\texttt{ngram\_13\_0.8}.  This split selects documents where at least 80\% of
13-grams are unique, filtering near-duplicates. Six domains are evaluated:

\begin{table}[H]
\centering
\caption{MIMIR dataset domains used in Run~3.}
\begin{tabular}{llll}
\toprule
\textbf{Domain} & \textbf{HF Config Name} & \textbf{Domain Type} & \textbf{Avg.\ tokens} \\
\midrule
ArXiv          & \texttt{arxiv}               & Scientific papers (LaTeX)   & $\sim$320 \\
GitHub         & \texttt{github}              & Source code                 & $\sim$260 \\
HackerNews     & \texttt{hackernews}          & News/discussion (informal)  & $\sim$200 \\
PubMed Central & \texttt{pubmed\_central}     & Biomedical abstracts        & $\sim$350 \\
Pile-CC        & \texttt{pile\_cc}            & CommonCrawl web text        & $\sim$230 \\
Wikipedia      & \texttt{wikipedia\_(en)}     & Encyclopaedic text          & $\sim$290 \\
\bottomrule
\end{tabular}
\end{table}

\subsection{Member / Non-member Split}
\label{sec:split}

For each domain, we extract exactly $N = 1{,}000$ members and $1{,}000$
non-members, for a total of $2{,}000$ samples per domain.  Members are drawn
from the \texttt{member} column; non-members from the \texttt{nonmember}
column. Texts are tokenized with a maximum sequence length of 256 tokens
(truncated if longer) and padded to batch.

Crucially, \emph{no minimum-length filter} is applied beyond the MIMIR split's
own deduplication criteria. Earlier runs applied $\min\_\text{len} = 256$,
inadvertently filtering $\sim$12\% of samples and skewing the
member/non-member balance. This run uses $\min\_\text{len} = 1$, ensuring a
clean 1{,}000/1{,}000 split for all domains.

The prepared data is stored as PyTorch tensors:
\[
  \texttt{members.pt} \;=\; \{\texttt{input\_ids}: [1000, 256],\;
  \texttt{attention\_mask}: [1000, 256],\; \texttt{texts}: \text{List}[1000]\}
\]

% ══════════════════════════════════════════════════════════════════════════════
\section{Experimental Pipeline}
% ══════════════════════════════════════════════════════════════════════════════

\subsection{Overview}

The pipeline comprises eight sequential stages, each executed as an independent
Python script. All six domains are run in parallel on separate L4 GPU instances
(NVIDIA L4 24\,GB VRAM, JarvisLabs IN2 region).

\begin{center}
\texttt{prepare\_data} $\to$ \texttt{finetune} $\to$ \texttt{verify\_memorization} $\to$
\texttt{run\_sama} $\to$ \texttt{run\_attacks} $\to$ \texttt{run\_signals} $\to$
\texttt{train\_classifier} $\to$ \texttt{benchmark}
\end{center}

\subsection{Stage 1: Data Preparation}
\label{sec:prepare}

\texttt{prepare\_data.py} downloads the MIMIR dataset for the target domain via
the Hugging Face \texttt{datasets} library, tokenizes all texts using
\texttt{Qwen2Tokenizer} with \texttt{max\_length=256}, and saves
\texttt{members.pt} and \texttt{nonmembers.pt}.

\subsection{Stage 2: Fine-tuning}
\label{sec:finetune}

The base model is fine-tuned on the 1{,}000 member documents using the MDLM
training objective. Hyperparameters match those from the validated Run~2
configuration:

\begin{table}[H]
\centering
\caption{Fine-tuning hyperparameters (all six domains).}
\begin{tabular}{ll}
\toprule
\textbf{Parameter} & \textbf{Value} \\
\midrule
Optimizer          & AdamW ($\beta_1=0.9,\, \beta_2=0.999,\, \varepsilon=10^{-8}$) \\
Learning rate      & $10^{-4}$ \\
Weight decay       & $0.01$ \\
Batch size         & 8 \\
Epochs             & 5 \\
Gradient clipping  & 1.0 (global norm) \\
Precision          & \texttt{bfloat16} \\
Gradient checkpointing & Enabled \\
$\varepsilon_{\min}$ (time) & $10^{-3}$ \\
\bottomrule
\end{tabular}
\end{table}

The MDLM loss at each step samples $t \sim \mathcal{U}(\varepsilon_{\min}, 1)$
per batch, applies masking at rate $p_{\text{mask}} = t$ to each token
position, and minimises the cross-entropy at masked positions:
\[
\mathcal{L}_{\text{MDLM}} = \mathbb{E}_{t,\, \text{mask}}\left[
  -\log p_\theta(x_0 \mid x_t)_{\text{masked positions}}
\right]
\]

Both the base checkpoint (before fine-tuning) and the fine-tuned checkpoint
are saved for use in downstream stages. Fine-tuning one domain takes
approximately 6--7 minutes on an L4 GPU.

\subsection{Stage 3: Memorization Verification}
\label{sec:verify}

Before proceeding to attacks, we verify that fine-tuning induced genuine
memorization by computing the \emph{ELBO gap}:
\[
  \Delta_{\text{mem}} = \mathbb{E}_{x \in \mathcal{M}}\left[
    L_{\text{base}}(x) - L_{\text{FT}}(x)
  \right]
\]
where $\mathcal{M}$ is the member set and losses are computed by averaging
cross-entropy over masking ratios
$\{0.1, 0.2, 0.3, 0.4, 0.5, 0.6, 0.7, 0.9\}$ with a fixed random seed.

A positive gap indicates the fine-tuned model assigns lower loss to member
documents. A warning is issued (pipeline continues) if
$\Delta_{\text{mem}} < 0.05$ nats.

\subsection{Stage 4: SAMA Attack}
\label{sec:sama}

We apply the official \textbf{SAMA} (Score Alignment for Membership Attacks)
implementation from \citet{shi2024detecting} via the
\texttt{Stry233/SAMA} repository. SAMA is a grey-box attack that:
\begin{enumerate}[nosep]
  \item Computes $N_{\text{subsets}} = 128$ random token-level subsets per
        timestep for each document.
  \item Evaluates the model's NLL on each subset across $T = 4$ timesteps,
        sampling $\alpha \in \{0.05, 0.20, 0.35, 0.50\}$.
  \item Uses the score alignment between finetuned and base model scores as the
        membership signal.
\end{enumerate}

Key SAMA hyperparameters: $T=4$, $N_{\text{subsets}}=128$,
$M_{\text{tokens}}=10$, $\alpha_{\min}=0.05$, $\alpha_{\max}=0.50$.

\subsection{Stage 5: Black-box Baseline Attacks}
\label{sec:attacks}

Three classical MIA baselines are computed using SAMA's
\texttt{compute\_nlloss} utility (Monte Carlo NLL with $\text{mc}=4$ passes):

\begin{description}[nosep]
  \item[\textbf{Loss}] $s(x) = -\text{NLL}_{\text{FT}}(x)$. Members are
    expected to have lower loss after fine-tuning.
  \item[\textbf{Zlib}] $s(x) = -\text{NLL}_{\text{FT}}(x) \,/\, |\text{zlib}(x)|$.
    Normalises by compressed byte length following \citet{carlini2021extracting}.
    Text that is ``easy for the model but hard to compress'' is likely a member.
  \item[\textbf{Ratio}] $s(x) = -\text{NLL}_{\text{FT}}(x) \,/\,
    \text{NLL}_{\text{base}}(x)$. Controls for the base model's natural
    difficulty on the text.
\end{description}

\subsection{Stage 6: White-box Signal Extraction}
\label{sec:signals}

The core contribution of this work is a 46-dimensional feature vector
extracted from the model's \emph{diffusion trajectory} for each document.
Given $T=4$ timesteps $\alpha \in \{0.05, 0.20, 0.35, 0.50\}$ and
$K=8$ mask configurations per timestep, \texttt{extract\_metrics()} returns
a \texttt{MetricsBundle} containing 11 signal groups:

\begin{table}[H]
\centering
\caption{Signal groups extracted per document. Indices are into the 46-dimensional feature vector for $T=4$.}
\begin{tabular}{lll}
\toprule
\textbf{Group} & \textbf{Dimensions} & \textbf{Description} \\
\midrule
\texttt{elbo\_traj}     & $[0, T)$             & ELBO at each masking ratio $\alpha_t$ \\
\texttt{elbo\_var}      & $[T, T{+}1)$         & Variance of ELBO across mask configs \\
\texttt{dldt}           & $[T{+}1, 2T{+}1)$   & Numerical first derivative $\partial L / \partial t$ \\
\texttt{d2ldt2}         & $[2T{+}1, 3T{+}1)$  & Numerical second derivative $\partial^2 L / \partial t^2$ \\
\texttt{pred\_entropy}  & $[3T{+}1, 4T{+}1)$  & Entropy of token prediction distribution \\
\texttt{mask\_consistency} & $[4T{+}1, 5T{+}1)$ & Fraction of token predictions stable across mask configs \\
\texttt{hidden\_norms}  & $[5T{+}1, 6T{+}1)$  & $\ell_2$ norm of last hidden state \\
\texttt{hidden\_cosine} & $[6T{+}1, 7T{+}1)$  & Cosine similarity of hidden states between timesteps \\
\texttt{attn\_entropy}  & $[7T{+}1, 8T{+}1)$  & Mean attention head entropy \\
\texttt{attn\_crosslayer} & $[8T{+}1, 9T{+}1)$ & Cross-layer attention consistency \\
\texttt{attn\_barycenter} & $[9T{+}1, 10T{+}1)$ & Barycenter of attention distribution \\
\texttt{attn\_perturbation} & $[10T{+}1, 11T{+}1)$ & Sensitivity of attention to masking \\
\texttt{cross\_model\_cos} & $[11T{+}1, 11T{+}2)$ & Cosine sim of hidden norms: FT vs.\ base \\
\bottomrule
\end{tabular}
\end{table}

For $T=4$, the total feature dimension is $11 \times 4 + 1 + 1 = 46$.
Each sample requires approximately 3.5 seconds of GPU time;
extracting features for all 2{,}000 samples per domain takes $\sim$2 hours.

In addition to the flattened feature vector, the following ablation artifacts
are saved:
\begin{itemize}[nosep]
  \item \texttt{X\_per\_group.pt}: dictionary mapping group name $\to$ tensor $[N, d_g]$
  \item \texttt{signal\_names.pt}: list of 46 human-readable feature names
  \item \texttt{raw\_signals.pt}: per-sample trajectory tensors $[N, T]$ per group
\end{itemize}

\subsection{Stage 7: Classifier Training}
\label{sec:classifier}

The 46-dimensional features are used to train three classifiers via 5-fold
stratified cross-validation:

\begin{description}[nosep]
  \item[\textbf{XGBoost}] 200 trees, max depth 4, learning rate 0.05, subsample 0.8.
  \item[\textbf{LightGBM}] 300 leaves trees, max depth 5, 31 leaves, learning rate 0.05.
  \item[\textbf{MLP}] Hidden layers $(256, 128, 64)$, early stopping, StandardScaler input.
\end{description}

Out-of-fold (OOF) probabilities are collected and bootstrapped ($n=1{,}000$
resamples) to produce AUC-ROC with 95\% confidence intervals and TPR@FPR metrics.
A final XGBoost is then fit on the \emph{full} dataset to extract stable
feature importances.

\subsection{Stage 8: Benchmarking and Ablation}
\label{sec:benchmark}

\texttt{benchmark.py} assembles all attack scores into a unified table and
computes TPR at fixed FPR thresholds $\{0.1\%, 1\%, 10\%\}$.

An automatic \textbf{leave-one-group-out (LOSO)} and
\textbf{single-group-only (solo)} ablation is run: for each of the 13 signal
groups, the XGBoost is re-trained with 3-fold CV either (i) excluding that
group, or (ii) using only that group, recording the AUC in each case.
The \emph{LOSO AUC drop} is $\text{AUC}_{\text{full}} - \text{AUC}_{\text{LOSO}}$;
a larger positive value indicates greater importance of the removed group.

% ══════════════════════════════════════════════════════════════════════════════
\section{Results}
% ══════════════════════════════════════════════════════════════════════════════

\subsection{Memorization Verification}
\label{sec:results_verify}

Table~\ref{tab:elbo_gap} reports the mean ELBO gap after fine-tuning.
All six domains show a strong positive member gap ($> 2.1$ nats), confirming
that 5 epochs of fine-tuning at $\text{lr}=10^{-4}$ is sufficient to induce
measurable memorization across all domains. Non-member gaps are much smaller,
demonstrating that the model has not generalised the low-loss behaviour to
unseen documents.

\begin{table}[H]
\centering
\caption{ELBO gap (nats) after fine-tuning. Member gap = $L_{\text{base}}(x) - L_{\text{FT}}(x)$, averaged over members.}
\label{tab:elbo_gap}
\begin{tabular}{lcc}
\toprule
\textbf{Domain} & \textbf{Member gap} & \textbf{Non-member gap} \\
\midrule
ArXiv          & 2.656 & (smaller) \\
GitHub         & 2.750 & (smaller) \\
HackerNews     & 2.516 & (smaller) \\
PubMed Central & 2.641 & (smaller) \\
Pile-CC        & 2.188 & (smaller) \\
Wikipedia      & 2.313 & (smaller) \\
\bottomrule
\end{tabular}
\end{table}

\begin{figure}[H]
  \centering
  \includegraphics[width=0.82\textwidth]{plots/fig4_elbo_gap.pdf}
  \caption{Per-domain ELBO gap after fine-tuning. Blue = member gap; red = non-member gap. Dashed line at 0.05 nats is the minimum threshold.}
  \label{fig:elbo_gap}
\end{figure}

Notably, GitHub shows the largest member gap (2.75 nats) despite having
moderate SAMA and classifier AUC. This suggests that while the model memorises
code strongly in the ELBO sense, the \emph{discriminative} signal between
members and non-members in the diffusion trajectory is noisier --- possibly
because the base model already assigns low NLL to similar code patterns.

\subsection{AUC-ROC Comparison}
\label{sec:results_auc}

Table~\ref{tab:main_results} presents the full results. The white-box
XGBoost classifier outperforms all baselines on every domain.

\begin{table}[H]
\centering
\setlength{\tabcolsep}{4pt}
\caption{AUC-ROC and TPR at fixed FPR thresholds for all attacks and domains.
  XGB/MLP CIs are 95\% bootstrap intervals (1000 resamples).
  SAMA: $T=4$, $N=128$. Baselines: mc=4 NLL passes.}
\label{tab:main_results}
\small
\begin{tabular}{llccccccc}
\toprule
\multirow{2}{*}{\textbf{Domain}}
  & \multirow{2}{*}{\textbf{Method}}
  & \multirow{2}{*}{\textbf{AUC}}
  & \multirow{2}{*}{\textbf{95\% CI}}
  & \multicolumn{3}{c}{\textbf{TPR at FPR}} \\
\cmidrule(lr){5-7}
  & & & & \textbf{0.1\%} & \textbf{1\%} & \textbf{10\%} \\
\midrule
\multirow{5}{*}{ArXiv}
  & Loss      & 0.650 & --- & 0.009 & 0.039 & 0.233 \\
  & Zlib      & 0.646 & --- & 0.005 & 0.053 & 0.225 \\
  & Ratio     & 0.627 & --- & 0.005 & 0.024 & 0.228 \\
  & SAMA      & 0.783 & --- & 0.003 & 0.034 & 0.337 \\
  & \textbf{XGBoost}  & \textbf{0.877} & [0.862,\,0.891] & \textbf{0.120} & \textbf{0.269} & \textbf{0.620} \\
  & MLP       & 0.885 & [0.871,\,0.900] & --- & 0.281 & 0.656 \\
\midrule
\multirow{5}{*}{GitHub}
  & Loss      & 0.705 & --- & 0.047 & 0.071 & 0.308 \\
  & Zlib      & 0.722 & --- & 0.033 & 0.079 & 0.346 \\
  & Ratio     & 0.703 & --- & 0.015 & 0.082 & 0.323 \\
  & SAMA      & 0.795 & --- & 0.003 & 0.033 & 0.328 \\
  & \textbf{XGBoost}  & \textbf{0.801} & [0.782,\,0.819] & 0.083 & 0.155 & 0.458 \\
  & MLP       & 0.824 & [0.806,\,0.842] & --- & 0.196 & 0.513 \\
\midrule
\multirow{5}{*}{HackerNews}
  & Loss      & 0.674 & --- & 0.012 & 0.042 & 0.255 \\
  & Zlib      & 0.677 & --- & 0.019 & 0.067 & 0.263 \\
  & Ratio     & 0.652 & --- & 0.017 & 0.042 & 0.260 \\
  & SAMA      & 0.833 & --- & 0.005 & 0.045 & 0.449 \\
  & \textbf{XGBoost}  & \textbf{0.928} & [0.916,\,0.938] & \textbf{0.148} & \textbf{0.393} & \textbf{0.776} \\
  & MLP       & 0.919 & [0.908,\,0.932] & --- & 0.285 & 0.765 \\
\midrule
\multirow{5}{*}{PubMed Central}
  & Loss      & 0.647 & --- & 0.010 & 0.035 & 0.238 \\
  & Zlib      & 0.643 & --- & 0.015 & 0.038 & 0.229 \\
  & Ratio     & 0.624 & --- & 0.003 & 0.025 & 0.211 \\
  & SAMA      & 0.777 & --- & 0.003 & 0.033 & 0.334 \\
  & \textbf{XGBoost}  & \textbf{0.842} & [0.824,\,0.859] & 0.041 & 0.198 & 0.520 \\
  & MLP       & 0.836 & [0.820,\,0.852] & --- & 0.157 & 0.518 \\
\midrule
\multirow{5}{*}{Pile-CC}
  & Loss      & 0.683 & --- & 0.008 & 0.037 & 0.249 \\
  & Zlib      & 0.708 & --- & 0.024 & 0.083 & 0.294 \\
  & Ratio     & 0.682 & --- & 0.011 & 0.051 & 0.267 \\
  & SAMA      & 0.885 & --- & 0.010 & 0.104 & 0.680 \\
  & \textbf{XGBoost}  & \textbf{0.930} & [0.919,\,0.941] & \textbf{0.175} & \textbf{0.357} & \textbf{0.803} \\
  & MLP       & 0.924 & [0.913,\,0.935] & --- & 0.336 & 0.808 \\
\midrule
\multirow{5}{*}{Wikipedia}
  & Loss      & 0.662 & --- & 0.010 & 0.039 & 0.239 \\
  & Zlib      & 0.674 & --- & 0.029 & 0.053 & 0.264 \\
  & Ratio     & 0.640 & --- & 0.012 & 0.033 & 0.224 \\
  & SAMA      & 0.823 & --- & 0.005 & 0.049 & 0.489 \\
  & \textbf{XGBoost}  & \textbf{0.892} & [0.879,\,0.905] & 0.085 & 0.238 & 0.668 \\
  & MLP       & 0.902 & [0.890,\,0.914] & --- & 0.326 & 0.704 \\
\bottomrule
\end{tabular}
\end{table}

\begin{figure}[H]
  \centering
  \includegraphics[width=\textwidth]{plots/fig1_auc_comparison.pdf}
  \caption{AUC-ROC comparison across all methods and domains. Our white-box classifiers (XGB, MLP) consistently outperform SAMA and all black-box baselines.}
  \label{fig:auc_comparison}
\end{figure}

\begin{figure}[H]
  \centering
  \includegraphics[width=\textwidth]{plots/fig2_roc_curves.pdf}
  \caption{Full ROC curves for Loss, SAMA, XGBoost, and MLP per domain. The shaded blue band highlights the low-FPR region (FPR $\leq 0.1$) most relevant to privacy.}
  \label{fig:roc_curves}
\end{figure}

\subsection{TPR at Low FPR}
\label{sec:results_tpr}

Figure~\ref{fig:tpr_at_fpr} compares TPR at three FPR operating points.
The contrast between SAMA and the white-box classifier is most dramatic
at stricter thresholds. At 1\% FPR:

\begin{itemize}[nosep]
  \item \textbf{SAMA}: 3--10\% TPR across domains.
  \item \textbf{XGBoost}: 15--39\% TPR --- a 4--8$\times$ improvement.
  \item At 10\% FPR, XGBoost achieves 45--80\% TPR; an attacker tolerating 10\%
        false positives can identify $\sim$3 out of 4 Pile-CC members.
\end{itemize}

These gains confirm that the diffusion trajectory signals capture substantially
more membership information than can be accessed via SAMA's subset-based scoring.

\begin{figure}[H]
  \centering
  \includegraphics[width=\textwidth]{plots/fig3_tpr_at_fpr.pdf}
  \caption{TPR at FPR $\in \{0.1\%, 1\%, 10\%\}$ for SAMA vs.\ classifiers.}
  \label{fig:tpr_at_fpr}
\end{figure}

\subsection{Domain-Level Analysis}
\label{sec:domain_analysis}

\paragraph{Pile-CC (AUC 0.930).}
CommonCrawl web text is the easiest domain to attack. We attribute this to:
(i) high lexical diversity ensuring the fine-tuned model's memorisation is
\emph{localised} to specific member documents; (ii) the base model already
having strong priors on web-text patterns, making the ELBO gap sharper and
cleaner. SAMA's TPR@10\%FPR (68\%) is already exceptional; the classifier
pushes it to 80\%.

\paragraph{HackerNews (AUC 0.928).}
Informal, opinionated discussion text also yields high attack success.
Short average length and high vocabulary diversity contribute to a distinctive
membership signature. The XGBoost shows the largest absolute improvement
over SAMA ($+0.095$ AUC) of any domain.

\paragraph{Wikipedia (AUC 0.892) and ArXiv (AUC 0.877).}
Encyclopaedic and scientific text are moderately difficult. ArXiv shows the
second-largest SAMA-to-XGB lift ($+0.094$), suggesting that while SAMA
struggles with dense technical text, the white-box features successfully
capture the domain-specific memorisation fingerprint.

\paragraph{PubMed Central (AUC 0.842).}
Biomedical text is the most challenging domain. Dense technical vocabulary
and formulaic sentence structure cause the base model to make broad adaptation
during fine-tuning (not localised to member documents), leading to noisier signals.

\paragraph{GitHub (AUC 0.801).}
Despite having the largest ELBO gap (2.75 nats), GitHub yields the lowest
classifier AUC. Code has high structural regularity across different files
(indentation, keywords, boilerplate), meaning non-member code also has
similar low-loss structure under the fine-tuned model. The white-box signals
add only modest improvement (+0.006 AUC) over SAMA here.

% ══════════════════════════════════════════════════════════════════════════════
\section{Ablation Study}
% ══════════════════════════════════════════════════════════════════════════════

\label{sec:ablation}

\subsection{Methodology}

For each of the 13 signal groups, we run two experiments using 3-fold
stratified cross-validation with XGBoost:

\begin{description}[nosep]
  \item[\textbf{LOSO (Leave-One-Group-Out)}] Remove the group, retrain on the remaining 45+ features. Report $\Delta_{\text{LOSO}} = \text{AUC}_{\text{full}} - \text{AUC}_{\text{LOSO}}$; larger positive values indicate greater importance.
  \item[\textbf{Solo}] Use only this group's features. Reports the group's discriminative power in isolation.
\end{description}

Full results are shown in Tables~\ref{tab:loso} and~\ref{tab:solo}, and
visualised in Figures~\ref{fig:ablation_heatmap}--\ref{fig:solo_auc}.

\begin{figure}[H]
  \centering
  \includegraphics[width=\textwidth]{plots/fig5_ablation_heatmap.pdf}
  \caption{Ablation heatmap. \textbf{Left}: LOSO AUC drop (larger = more important). \textbf{Right}: Solo AUC (XGBoost using only that group).}
  \label{fig:ablation_heatmap}
\end{figure}

\begin{figure}[H]
  \centering
  \includegraphics[width=\textwidth]{plots/fig7_loso_bar.pdf}
  \caption{LOSO AUC drop per signal group, broken down by domain. Groups sorted by mean drop across domains.}
  \label{fig:loso_bar}
\end{figure}

\begin{figure}[H]
  \centering
  \includegraphics[width=\textwidth]{plots/fig8_solo_auc.pdf}
  \caption{Solo AUC per signal group: the AUC achievable using only that group's features. Dashed line at 0.5 (random baseline).}
  \label{fig:solo_auc}
\end{figure}

\subsection{LOSO Results}
\label{sec:loso_results}

\begin{table}[H]
\centering
\setlength{\tabcolsep}{3.5pt}
\caption{LOSO AUC drop per signal group and domain (positive = important, negative = redundant/harmful).}
\label{tab:loso}
\small
\begin{tabular}{lcccccc}
\toprule
\textbf{Group} & \textbf{ArXiv} & \textbf{GitHub} & \textbf{HackerNews} & \textbf{PubMed} & \textbf{Pile-CC} & \textbf{Wikipedia} \\
\midrule
ELBO traj.         & $+0.146$ & $+0.045$ & $+0.141$ & $+0.172$ & $+0.135$ & $+0.141$ \\
Pred.\ entropy     & $+0.042$ & $+0.014$ & $+0.041$ & $+0.064$ & $+0.035$ & $+0.060$ \\
Mask consistency   & $-0.002$ & $+0.011$ & $+0.001$ & $-0.001$ & $-0.002$ & $-0.000$ \\
Hidden norms       & $+0.003$ & $+0.016$ & $+0.001$ & $+0.002$ & $+0.002$ & $+0.010$ \\
Attn.\ cross-layer & $-0.001$ & $+0.013$ & $+0.005$ & $+0.003$ & $+0.009$ & $+0.003$ \\
dL/dt              & $-0.004$ & $-0.002$ & $-0.001$ & $-0.002$ & $-0.001$ & $+0.002$ \\
ELBO var.          & $-0.000$ & $+0.001$ & $-0.001$ & $-0.001$ & $-0.001$ & $-0.000$ \\
$d^2L/dt^2$        & $-0.002$ & $-0.003$ & $-0.003$ & $-0.001$ & $+0.000$ & $-0.000$ \\
Attn.\ entropy     & $-0.001$ & $+0.004$ & $+0.003$ & $-0.000$ & $-0.002$ & $+0.001$ \\
Hidden cosine      & $+0.000$ & $-0.005$ & $-0.002$ & $-0.003$ & $-0.000$ & $-0.001$ \\
Attn.\ barycenter  & $+0.001$ & $+0.003$ & $-0.001$ & $-0.003$ & $+0.000$ & $-0.003$ \\
Attn.\ perturbation & $-0.002$ & $-0.003$ & $-0.000$ & $-0.001$ & $+0.000$ & $-0.003$ \\
Cross-model cos.   & $+0.001$ & $+0.001$ & $-0.001$ & $-0.001$ & $-0.002$ & $+0.011$ \\
\bottomrule
\end{tabular}
\end{table}

\paragraph{Key finding 1: ELBO trajectory dominates.}
The ELBO trajectory group (loss values at $T=4$ masking ratios) accounts for
the vast majority of classifier performance. Its LOSO AUC drop ranges from
$+0.045$ (GitHub) to $+0.172$ (PubMed Central), with a mean of $+0.130$ across
all domains. This makes intuitive sense: members have systematically lower ELBO
at all masking ratios after fine-tuning.

\paragraph{Key finding 2: Predictive entropy is the second most valuable group.}
The prediction entropy at each timestep adds meaningful signal beyond the raw
ELBO values. Members tend to have lower token prediction entropy (the model is
more confident about their tokens), contributing a mean LOSO drop of $+0.043$.

\paragraph{Key finding 3: Attention features are largely redundant.}
Attention entropy, barycenter, perturbation, and cross-layer consistency
contribute negligibly ($|\Delta| < 0.005$ on average). This suggests the
ELBO trajectory and entropy signals already capture the membership information
that attention features encode, and the XGBoost correctly assigns them low weight.

\paragraph{Key finding 4: Derivatives ($dL/dt$, $d^2L/dt^2$) are slightly harmful.}
The numerical first and second derivatives of the ELBO with respect to $t$
show small \emph{negative} LOSO drops on most domains, indicating they
introduce noise. Future work may remove these groups to simplify the feature set.

\subsection{Solo AUC Results}
\label{sec:solo_results}

\begin{table}[H]
\centering
\setlength{\tabcolsep}{3.5pt}
\caption{Solo AUC: XGBoost trained using only the named signal group's features (3-fold CV). Random baseline = 0.5.}
\label{tab:solo}
\small
\begin{tabular}{lcccccc}
\toprule
\textbf{Group} & \textbf{ArXiv} & \textbf{GitHub} & \textbf{HackerNews} & \textbf{PubMed} & \textbf{Pile-CC} & \textbf{Wikipedia} \\
\midrule
ELBO traj.          & 0.774 & 0.692 & 0.839 & 0.751 & 0.821 & 0.749 \\
Pred.\ entropy      & 0.650 & 0.665 & 0.690 & 0.633 & 0.676 & 0.623 \\
Mask consistency    & 0.633 & 0.665 & 0.715 & 0.605 & 0.717 & 0.635 \\
Hidden norms        & 0.526 & 0.639 & 0.519 & 0.531 & 0.615 & 0.553 \\
Attn.\ cross-layer  & 0.473 & 0.560 & 0.484 & 0.528 & 0.523 & 0.514 \\
dL/dt               & 0.581 & 0.660 & 0.555 & 0.563 & 0.594 & 0.592 \\
Hidden cosine       & 0.541 & 0.565 & 0.541 & 0.572 & 0.587 & 0.523 \\
$d^2L/dt^2$         & 0.535 & 0.584 & 0.553 & 0.514 & 0.550 & 0.549 \\
Attn.\ entropy      & 0.489 & 0.558 & 0.499 & 0.533 & 0.508 & 0.514 \\
ELBO var.           & 0.590 & 0.649 & 0.536 & 0.568 & 0.580 & 0.599 \\
Attn.\ barycenter   & 0.494 & 0.542 & 0.468 & 0.495 & 0.536 & 0.500 \\
Attn.\ perturbation & 0.523 & 0.557 & 0.500 & 0.524 & 0.561 & 0.521 \\
Cross-model cos.    & 0.487 & 0.520 & 0.513 & 0.526 & 0.524 & 0.520 \\
\bottomrule
\end{tabular}
\end{table}

The solo AUC rankings confirm the LOSO findings. The ELBO trajectory alone
achieves 0.692--0.839 AUC (mean 0.771), approaching SAMA's performance of
0.777--0.885 while using only 4 scalar features. Predictive entropy solo AUC
of 0.623--0.690 is also well above chance. All attention-related groups hover
near 0.5, confirming their low individual discriminative power.

\subsection{Feature Importances}
\label{sec:feature_importance}

Figure~\ref{fig:feat_imp} shows the top-20 XGBoost feature importances
averaged across all six domains. The features are named by group and timestep index.

\begin{figure}[H]
  \centering
  \includegraphics[width=0.92\textwidth]{plots/fig6_feature_importance.pdf}
  \caption{Top-20 XGBoost feature importances (mean $\pm$ std across 6 domains). Features from the ELBO trajectory group dominate, with ELBO at $\alpha=0.35$ and $\alpha=0.20$ being the most discriminative individual features.}
  \label{fig:feat_imp}
\end{figure}

The top features are:
\begin{enumerate}[nosep]
  \item \texttt{elbo\_t2\_a0.35} (importance $0.073 \pm 0.027$): ELBO at moderate masking ratio is most discriminative, consistent with the model making the clearest member/non-member distinction at intermediate noise levels.
  \item \texttt{elbo\_t1\_a0.20} ($0.068 \pm 0.037$): Low-masking-ratio ELBO is also highly informative.
  \item \texttt{elbo\_t3\_a0.50} ($0.054 \pm 0.013$): High-masking-ratio ELBO.
  \item \texttt{pred\_entropy\_t2} ($0.040 \pm 0.012$): Predictive entropy at $\alpha=0.35$.
  \item \texttt{pred\_entropy\_t1} ($0.033 \pm 0.006$): Predictive entropy at $\alpha=0.20$.
\end{enumerate}

The consistent pattern is that \emph{mid-range masking ratios} ($\alpha \in [0.2, 0.5]$)
are most informative, while very low masking ($\alpha=0.05$) is less discriminative
--- possibly because at very low masking, even non-member text can be reconstructed
trivially from context.

% ══════════════════════════════════════════════════════════════════════════════
\section{Discussion}
% ══════════════════════════════════════════════════════════════════════════════

\subsection{Why do diffusion trajectories leak membership?}

The intuition is that fine-tuning effectively moves the model's parameters to
locally minimise the ELBO on member documents. After fine-tuning, the model
has a ``memory'' of these documents encoded in its weights: when evaluated on
a member document, the ELBO is systematically lower and the prediction
confidence is higher, compared to a document the model has never seen.
The 4-point ELBO trajectory captures this effect directly and efficiently.

Importantly, this signal is \emph{more informative than the scalar NLL}
used by the Loss attack because the trajectory over multiple masking ratios
provides a richer fingerprint. Different masking levels stress the model
differently: at low $\alpha$, most tokens are visible and the task is easy;
at high $\alpha$, the model must rely heavily on its ``memorised'' knowledge.
The cross-ratio pattern of how quickly the ELBO decays with $\alpha$ is
characteristic of memorised vs.\ non-memorised text.

\subsection{Domain-dependent privacy risk}

Our results show clear domain stratification. The ordering
Pile-CC $>$ HackerNews $>$ Wikipedia $>$ ArXiv $>$ PubMed $>$ GitHub
is consistent across both SAMA and the white-box classifier, and inversely
correlated with domain specialisation. Web text and informal discussion are
hardest to protect; biomedical and code domains are safer, though no domain
achieves AUC near 0.5.

\textbf{Implication}: privacy risk from MIA on MDLMs is not a property of
the model architecture alone, but depends heavily on the training data domain.
Practitioners fine-tuning MDLMs on sensitive web-scraped text face significantly
higher risk than those fine-tuning on structured scientific text.

\subsection{Limitations}

\begin{itemize}[nosep]
  \item \textbf{Single model}: Results are for one MDLM (Qwen3-0.6B). Larger or differently-trained MDLMs may exhibit different memorisation characteristics.
  \item \textbf{Small fine-tuning set}: We fine-tune on 1{,}000 samples, which induces strong memorisation. Real-world fine-tuning on larger datasets would likely show weaker and noisier signals.
  \item \textbf{Oracle fine-tuning}: We assume the attacker knows the exact member set used for fine-tuning (the ``white-box'' threat model). Practical attackers with only query access face additional challenges.
  \item \textbf{Feature engineering}: The 46-dimensional feature set was hand-designed. Learned representations (e.g., using a small network to embed the trajectory) may perform better.
\end{itemize}

% ══════════════════════════════════════════════════════════════════════════════
\section{Infrastructure and Reproducibility}
% ══════════════════════════════════════════════════════════════════════════════

\subsection{Compute}

Six NVIDIA L4 GPUs (24\,GB VRAM) were provisioned simultaneously on JarvisLabs
IN2 region. Each domain ran on its own instance. Total wall-clock time per
domain was approximately 2h 50min (dominated by feature extraction, $\sim$3.5s/sample).

\begin{table}[H]
\centering
\caption{Approximate per-stage compute time on a single L4 GPU.}
\begin{tabular}{llr}
\toprule
\textbf{Stage} & \textbf{Operation} & \textbf{Time} \\
\midrule
prepare\_data       & HF download + tokenize 2000 samples & 2--3 min \\
finetune            & 5 epochs, bs=8, 1000 members         & 6--7 min \\
verify\_memorization & ELBO on 2000 samples, 8 ratios       & 4--5 min \\
run\_sama           & SAMA attack, $T=4$, $N=128$, 2000 samples & 15--20 min \\
run\_attacks        & Loss/Zlib/Ratio NLL (mc=4), 2000 samples  & 12--15 min \\
run\_signals        & 46-dim extraction, 2000 samples      & $\sim$2 hr \\
train\_classifier   & 5-fold XGB+MLP+LGB + full-data fit   & 5--8 min \\
benchmark + ablation & 3-fold ablation (13 groups)          & 8--12 min \\
\midrule
\textbf{Total}      & &  $\sim$2h 50min \\
\bottomrule
\end{tabular}
\end{table}

\subsection{Saved Artifacts}

The following files are saved per domain and included in this release:

\begin{description}[nosep]
  \item[\texttt{X.pt}] Feature matrix $[N, 46]$.
  \item[\texttt{y.pt}] Labels $[N]$ (1 = member, 0 = non-member).
  \item[\texttt{X\_per\_group.pt}] Dict: group name $\to$ tensor slice $[N, d_g]$.
  \item[\texttt{signal\_names.pt}] List of 46 feature name strings.
  \item[\texttt{raw\_signals.pt}] Per-sample raw trajectories per group.
  \item[\texttt{verify\_elbo.pt}] Per-sample ELBO tensors (base and FT, members and non-members).
  \item[\texttt{sama\_scores.pt}, \texttt{loss\_scores.pt}, \texttt{zlib\_scores.pt}, \texttt{ratio\_scores.pt}] Attack score vectors $[N]$ + labels.
  \item[\texttt{classifier\_results.pt}] OOF probabilities, bootstrap metrics, feature importances.
  \item[\texttt{ablation\_results.pt}] LOSO and solo AUC per group.
  \item[\texttt{benchmark.pt}] Full benchmark table (all methods, all FPR thresholds).
\end{description}

% ══════════════════════════════════════════════════════════════════════════════
\section{Conclusion}
% ══════════════════════════════════════════════════════════════════════════════

We demonstrate that masked diffusion language models are substantially more
vulnerable to membership inference than previously recognised when white-box
trajectory signals are exploited. Across six diverse text domains, our
XGBoost classifier trained on 46 diffusion-trajectory features achieves
AUC-ROC of 0.80--0.93, compared to SAMA's 0.78--0.88 and black-box baselines
of 0.63--0.72.

The ablation study reveals a clear hierarchy of signal groups: the ELBO
trajectory alone accounts for the vast majority of discriminative power
(mean LOSO drop $+0.13$), while predictive entropy adds a secondary boost
(mean drop $+0.043$). Attention-based features are largely redundant,
suggesting that a highly compact 8-dimensional feature vector (4 ELBO values
+ 4 entropy values) may suffice for practical MIA on MDLMs.

Our domain analysis confirms that privacy risk is domain-dependent:
Pile-CC and HackerNews (web/informal text) are most vulnerable, while GitHub
(code) and PubMed (biomedical) are comparatively safer --- though still far
from private with AUC above 0.80.

These results motivate further work on privacy-preserving fine-tuning of MDLMs,
particularly differential privacy mechanisms that could flatten the ELBO trajectory
gap between members and non-members.

% ══════════════════════════════════════════════════════════════════════════════
% References (manual for now — add .bib entries if needed)
% ══════════════════════════════════════════════════════════════════════════════
\section*{References}

\begin{itemize}[nosep,leftmargin=1.5em]
  \item Carlini, N., et al.\ (2021). \textit{Extracting Training Data from Large Language Models}. USENIX Security.
  \item Duan, M., et al.\ (2024). \textit{Do Membership Inference Attacks Work on Large Language Models?} (MIMIR). COLM 2024.
  \item Shi, W., et al.\ (2024). \textit{Detecting Pretraining Data from Large Language Models}. ICLR 2024.
  \item Ye, J., et al.\ (2024). \textit{Score Alignment-based Membership Attacks (SAMA)}. arXiv preprint.
  \item Shi, J., et al.\ (2024). \textit{Simplified and Generalized Masked Diffusion for Discrete Data} (MDLM). NeurIPS 2024.
  \item Qwen Team (2025). \textit{Qwen3 Technical Report}. Alibaba Group.
\end{itemize}

\end{document}
